% --- Archon physics formalization source begin ---
% archon:physics
% archon:covers PhyXMiniProblems/problem_phyx_mini_0777.lean
% archon:source-report reports/phyx_mini/problem_phyx_mini_0777.source.json

\chapter{Physics problem phyx\_mini\_0777}
\label{ch:PhyXMiniProblems_problem_phyx_mini_0777}

\paragraph{Problem source.}
\#\# Physical scenario

A square metal plate \textbackslash{}(0.180\textbackslash{}text\{ m\}\textbackslash{}) on each side is pivoted about an axis through point \textbackslash{}(O\textbackslash{}) at its center and perpendicular to the plate as shown in figure.

\#\# Question

Calculate the net torque about this axis due to the three forces shown in the figure if the magnitudes of the forces are \textbackslash{}(F\_1 = 18.0\textbackslash{}text\{ N\}\textbackslash{}), \textbackslash{}(F\_2 = 26.0\textbackslash{}text\{ N\}\textbackslash{}), and \textbackslash{}(F\_3 = 14.0\textbackslash{}text\{ N\}\textbackslash{}). The plate and all forces are in the plane of the page.

\#\# Answer choices

- A: A: \textbackslash{}(1.50 \textbackslash{}text\{ N\} \textbackslash{}cdot \textbackslash{}text\{m\} \textbackslash{})

- B: B: \textbackslash{}(2.50 \textbackslash{}text\{ N\} \textbackslash{}cdot \textbackslash{}text\{m\} \textbackslash{})

- C: C: \textbackslash{}(5.00 \textbackslash{}text\{ N\} \textbackslash{}cdot \textbackslash{}text\{m\} \textbackslash{})

- D: D: \textbackslash{}(4.50 \textbackslash{}text\{ N\} \textbackslash{}cdot \textbackslash{}text\{m\} \textbackslash{})

\#\# Figure caption (auxiliary; use the image as primary evidence)

The image depicts a square object with a side length of 0.180 m. The square is titled, and a point labeled "O" is marked at the center. Three forces are applied to the square:

1. Force \textbackslash{}( F\_1 \textbackslash{}), directed vertically downward along the right side.
2. Force \textbackslash{}( F\_2 \textbackslash{}), directed vertically upward along the left side.
3. Force \textbackslash{}( F\_3 \textbackslash{}), applied at a 45-degree angle from the bottom right corner, directed away from the square.

The sides of the square are labeled with their length, 0.180 m. A small arrow indicates the 45-degree angle at which \textbackslash{}( F\_3 \textbackslash{}) is applied. The forces \textbackslash{}( F\_1 \textbackslash{}), \textbackslash{}( F\_2 \textbackslash{}), and \textbackslash{}( F\_3 \textbackslash{}) are represented by red arrows.

\#\# Dataset metadata

Rotational Motion; Spatial Relation Reasoning, Multi-Formula Reasoning

\paragraph{Recorded answer/context.}
B: B: \textbackslash{}(2.50 \textbackslash{}text\{ N\} \textbackslash{}cdot \textbackslash{}text\{m\} \textbackslash{})

\paragraph{Figure/image path.}
/root/proposal\_for\_physic/science-mango/phyx\_mini\_run/phyx\_data/test\_image/777.png

\paragraph{Formalization target.}
create a compiling Lean file with sorry bodies at `PhyXMiniProblems/problem\_phyx\_mini\_0777.lean`. The Lean declarations must preserve the physical quantities, dimensions or dimensional roles, figure labels, governing-law hypotheses, and final relation expressed by this problem.
Use Mathlib/Physlib names found through LeanExplore where available. If a physics API is missing, introduce faithful local abstractions rather than scalar placeholder aliases.

\begin{theorem}[Physics formalization target]
\label{thm:physics:phyx_mini_0777:target}
\lean{PhyXMiniProblems.ProblemPhyXMini0777.problem_phyx_mini_0777}
\uses{def:physics:phyx-mini-0777:phyxminiproblems-problemphyxmini0777-torqueinnewtonmeters, def:physics:phyx-mini-0777:phyxminiproblems-problemphyxmini0777-squareplatetorquesetup, def:physics:phyx-mini-0777:phyxminiproblems-problemphyxmini0777-matchesproblemstatement, def:physics:phyx-mini-0777:phyxminiproblems-problemphyxmini0777-matchesprimaryfigure, def:physics:phyx-mini-0777:phyxminiproblems-problemphyxmini0777-matchessquareplategeometry, def:physics:phyx-mini-0777:phyxminiproblems-problemphyxmini0777-matchesdisplayedforcegeometry, def:physics:phyx-mini-0777:phyxminiproblems-problemphyxmini0777-satisfiesplanartorquelaw, def:physics:phyx-mini-0777:phyxminiproblems-problemphyxmini0777-recordedanswerchoice, def:physics:phyx-mini-0777:phyxminiproblems-problemphyxmini0777-roundstodisplayedtorque}
The assigned autoformalize agent should translate this physics problem into Lean declarations in the covered file, with theorem and lemma proof bodies written as `by sorry`.
\end{theorem}
\begin{proof}
This is an autoformalization task, not a proof task. Produce statements that can later be proved by the physics prover without weakening the physical model.
\end{proof}

% --- Archon declaration topology begin ---
\section{Declaration topology}
\begin{definition}[force Dimension]
  \label{def:physics:phyx-mini-0777:phyxminiproblems-problemphyxmini0777-forcedimension}
  \lean{PhyXMiniProblems.ProblemPhyXMini0777.forceDimension}
  The physical dimension M L T⁻² of force
\end{definition}
\begin{proof}
  This is the declared model definition of force Dimension.
\end{proof}
\begin{definition}[torque Dimension]
  \label{def:physics:phyx-mini-0777:phyxminiproblems-problemphyxmini0777-torquedimension}
  \lean{PhyXMiniProblems.ProblemPhyXMini0777.torqueDimension}
  \uses{def:physics:phyx-mini-0777:phyxminiproblems-problemphyxmini0777-forcedimension}
  The physical dimension M L² T⁻² of torque
\end{definition}
\begin{proof}
  This is the declared model definition of torque Dimension.
\end{proof}
\begin{definition}[Planar Vector]
  \label{def:physics:phyx-mini-0777:phyxminiproblems-problemphyxmini0777-planarvector}
  \lean{PhyXMiniProblems.ProblemPhyXMini0777.PlanarVector}
  Vectors in the plane of the plate
\end{definition}
\begin{proof}
  This is the declared model definition of Planar Vector.
\end{proof}
\begin{definition}[Length Quantity]
  \label{def:physics:phyx-mini-0777:phyxminiproblems-problemphyxmini0777-lengthquantity}
  \lean{PhyXMiniProblems.ProblemPhyXMini0777.LengthQuantity}
  A unit-independent nonnegative physical length
\end{definition}
\begin{proof}
  This is the declared model definition of Length Quantity.
\end{proof}
\begin{definition}[Planar Position Quantity]
  \label{def:physics:phyx-mini-0777:phyxminiproblems-problemphyxmini0777-planarpositionquantity}
  \lean{PhyXMiniProblems.ProblemPhyXMini0777.PlanarPositionQuantity}
  \uses{def:physics:phyx-mini-0777:phyxminiproblems-problemphyxmini0777-planarvector}
  A unit-independent position vector in the plane of the plate
\end{definition}
\begin{proof}
  This is the declared model definition of Planar Position Quantity.
\end{proof}
\begin{definition}[Planar Force Quantity]
  \label{def:physics:phyx-mini-0777:phyxminiproblems-problemphyxmini0777-planarforcequantity}
  \lean{PhyXMiniProblems.ProblemPhyXMini0777.PlanarForceQuantity}
  \uses{def:physics:phyx-mini-0777:phyxminiproblems-problemphyxmini0777-forcedimension, def:physics:phyx-mini-0777:phyxminiproblems-problemphyxmini0777-planarvector}
  A unit-independent force vector in the plane of the plate
\end{definition}
\begin{proof}
  This is the declared model definition of Planar Force Quantity.
\end{proof}
\begin{definition}[Signed Torque Quantity]
  \label{def:physics:phyx-mini-0777:phyxminiproblems-problemphyxmini0777-signedtorquequantity}
  \lean{PhyXMiniProblems.ProblemPhyXMini0777.SignedTorqueQuantity}
  \uses{def:physics:phyx-mini-0777:phyxminiproblems-problemphyxmini0777-torquedimension}
  A signed torque component about the axis perpendicular to the plate
\end{definition}
\begin{proof}
  This is the declared model definition of Signed Torque Quantity.
\end{proof}
\begin{definition}[x Axis]
  \label{def:physics:phyx-mini-0777:phyxminiproblems-problemphyxmini0777-xaxis}
  \lean{PhyXMiniProblems.ProblemPhyXMini0777.xAxis}
  Coordinate 0, pointing right in the primary image
\end{definition}
\begin{proof}
  This is the declared model definition of x Axis.
\end{proof}
\begin{definition}[y Axis]
  \label{def:physics:phyx-mini-0777:phyxminiproblems-problemphyxmini0777-yaxis}
  \lean{PhyXMiniProblems.ProblemPhyXMini0777.yAxis}
  Coordinate 1, pointing upward in the primary image
\end{definition}
\begin{proof}
  This is the declared model definition of y Axis.
\end{proof}
\begin{definition}[nonnegative Readout]
  \label{def:physics:phyx-mini-0777:phyxminiproblems-problemphyxmini0777-nonnegativereadout}
  \lean{PhyXMiniProblems.ProblemPhyXMini0777.nonnegativeReadout}
  Read a nonnegative physical scalar in a coherent unit system
\end{definition}
\begin{proof}
  This is the declared model definition of nonnegative Readout.
\end{proof}
\begin{definition}[planar Vector Readout]
  \label{def:physics:phyx-mini-0777:phyxminiproblems-problemphyxmini0777-planarvectorreadout}
  \lean{PhyXMiniProblems.ProblemPhyXMini0777.planarVectorReadout}
  \uses{def:physics:phyx-mini-0777:phyxminiproblems-problemphyxmini0777-planarvector}
  Read a planar physical vector in a coherent unit system
\end{definition}
\begin{proof}
  This is the declared model definition of planar Vector Readout.
\end{proof}
\begin{definition}[signed Scalar Readout]
  \label{def:physics:phyx-mini-0777:phyxminiproblems-problemphyxmini0777-signedscalarreadout}
  \lean{PhyXMiniProblems.ProblemPhyXMini0777.signedScalarReadout}
  Read a signed physical scalar in a coherent unit system
\end{definition}
\begin{proof}
  This is the declared model definition of signed Scalar Readout.
\end{proof}
\begin{definition}[length In Meters]
  \label{def:physics:phyx-mini-0777:phyxminiproblems-problemphyxmini0777-lengthinmeters}
  \lean{PhyXMiniProblems.ProblemPhyXMini0777.lengthInMeters}
  \uses{def:physics:phyx-mini-0777:phyxminiproblems-problemphyxmini0777-lengthquantity, def:physics:phyx-mini-0777:phyxminiproblems-problemphyxmini0777-nonnegativereadout}
  Metre readout of a nonnegative physical length
\end{definition}
\begin{proof}
  This is the declared model definition of length In Meters.
\end{proof}
\begin{definition}[position In Meters]
  \label{def:physics:phyx-mini-0777:phyxminiproblems-problemphyxmini0777-positioninmeters}
  \lean{PhyXMiniProblems.ProblemPhyXMini0777.positionInMeters}
  \uses{def:physics:phyx-mini-0777:phyxminiproblems-problemphyxmini0777-planarvector, def:physics:phyx-mini-0777:phyxminiproblems-problemphyxmini0777-planarpositionquantity, def:physics:phyx-mini-0777:phyxminiproblems-problemphyxmini0777-planarvectorreadout}
  Cartesian metre readout of a position in the plate
\end{definition}
\begin{proof}
  This is the declared model definition of position In Meters.
\end{proof}
\begin{definition}[force Vector In Newtons]
  \label{def:physics:phyx-mini-0777:phyxminiproblems-problemphyxmini0777-forcevectorinnewtons}
  \lean{PhyXMiniProblems.ProblemPhyXMini0777.forceVectorInNewtons}
  \uses{def:physics:phyx-mini-0777:phyxminiproblems-problemphyxmini0777-planarvector, def:physics:phyx-mini-0777:phyxminiproblems-problemphyxmini0777-planarforcequantity, def:physics:phyx-mini-0777:phyxminiproblems-problemphyxmini0777-planarvectorreadout}
  Cartesian newton readout of a planar force
\end{definition}
\begin{proof}
  This is the declared model definition of force Vector In Newtons.
\end{proof}
\begin{definition}[force Magnitude In Newtons]
  \label{def:physics:phyx-mini-0777:phyxminiproblems-problemphyxmini0777-forcemagnitudeinnewtons}
  \lean{PhyXMiniProblems.ProblemPhyXMini0777.forceMagnitudeInNewtons}
  \uses{def:physics:phyx-mini-0777:phyxminiproblems-problemphyxmini0777-planarforcequantity, def:physics:phyx-mini-0777:phyxminiproblems-problemphyxmini0777-forcevectorinnewtons}
  Euclidean magnitude, in newtons, of a planar physical force
\end{definition}
\begin{proof}
  This is the declared model definition of force Magnitude In Newtons.
\end{proof}
\begin{definition}[torque In Newton Meters]
  \label{def:physics:phyx-mini-0777:phyxminiproblems-problemphyxmini0777-torqueinnewtonmeters}
  \lean{PhyXMiniProblems.ProblemPhyXMini0777.torqueInNewtonMeters}
  \uses{def:physics:phyx-mini-0777:phyxminiproblems-problemphyxmini0777-signedtorquequantity, def:physics:phyx-mini-0777:phyxminiproblems-problemphyxmini0777-signedscalarreadout}
  Signed newton-metre readout, positive for counterclockwise torque
\end{definition}
\begin{proof}
  This is the declared model definition of torque In Newton Meters.
\end{proof}
\begin{definition}[planar Moment Readout]
  \label{def:physics:phyx-mini-0777:phyxminiproblems-problemphyxmini0777-planarmomentreadout}
  \lean{PhyXMiniProblems.ProblemPhyXMini0777.planarMomentReadout}
  \uses{def:physics:phyx-mini-0777:phyxminiproblems-problemphyxmini0777-planarvector, def:physics:phyx-mini-0777:phyxminiproblems-problemphyxmini0777-xaxis, def:physics:phyx-mini-0777:phyxminiproblems-problemphyxmini0777-yaxis}
  The out-of-plane component of the moment of a force whose lever-arm and force coordinates have been read in one coherent unit system
\end{definition}
\begin{proof}
  This is the declared model definition of planar Moment Readout.
\end{proof}
\begin{definition}[Plate Point]
  \label{def:physics:phyx-mini-0777:phyxminiproblems-problemphyxmini0777-platepoint}
  \lean{PhyXMiniProblems.ProblemPhyXMini0777.PlatePoint}
  Distinguished center and corners of the square plate
\end{definition}
\begin{proof}
  This is the declared model definition of Plate Point.
\end{proof}
\begin{definition}[Applied Force Label]
  \label{def:physics:phyx-mini-0777:phyxminiproblems-problemphyxmini0777-appliedforcelabel}
  \lean{PhyXMiniProblems.ProblemPhyXMini0777.AppliedForceLabel}
  The three force labels printed beside the red arrows
\end{definition}
\begin{proof}
  This is the declared model definition of Applied Force Label.
\end{proof}
\begin{definition}[Side Length Label]
  \label{def:physics:phyx-mini-0777:phyxminiproblems-problemphyxmini0777-sidelengthlabel}
  \lean{PhyXMiniProblems.ProblemPhyXMini0777.SideLengthLabel}
  The two separate side-length labels visible in the primary image
\end{definition}
\begin{proof}
  This is the declared model definition of Side Length Label.
\end{proof}
\begin{definition}[Arrow Direction]
  \label{def:physics:phyx-mini-0777:phyxminiproblems-problemphyxmini0777-arrowdirection}
  \lean{PhyXMiniProblems.ProblemPhyXMini0777.ArrowDirection}
  Qualitative directions of the force arrows in the image
\end{definition}
\begin{proof}
  This is the declared model definition of Arrow Direction.
\end{proof}
\begin{definition}[Plate Material]
  \label{def:physics:phyx-mini-0777:phyxminiproblems-problemphyxmini0777-platematerial}
  \lean{PhyXMiniProblems.ProblemPhyXMini0777.PlateMaterial}
  Material category stated for the plate
\end{definition}
\begin{proof}
  This is the declared model definition of Plate Material.
\end{proof}
\begin{definition}[expected Application Point]
  \label{def:physics:phyx-mini-0777:phyxminiproblems-problemphyxmini0777-expectedapplicationpoint}
  \lean{PhyXMiniProblems.ProblemPhyXMini0777.expectedApplicationPoint}
  \uses{def:physics:phyx-mini-0777:phyxminiproblems-problemphyxmini0777-platepoint, def:physics:phyx-mini-0777:phyxminiproblems-problemphyxmini0777-appliedforcelabel}
  The corner at which each force arrow is applied in image 777.png
\end{definition}
\begin{proof}
  This is the declared model definition of expected Application Point.
\end{proof}
\begin{definition}[expected Arrow Direction]
  \label{def:physics:phyx-mini-0777:phyxminiproblems-problemphyxmini0777-expectedarrowdirection}
  \lean{PhyXMiniProblems.ProblemPhyXMini0777.expectedArrowDirection}
  \uses{def:physics:phyx-mini-0777:phyxminiproblems-problemphyxmini0777-appliedforcelabel, def:physics:phyx-mini-0777:phyxminiproblems-problemphyxmini0777-arrowdirection}
  The arrow direction read directly from image 777.png
\end{definition}
\begin{proof}
  This is the declared model definition of expected Arrow Direction.
\end{proof}
\begin{definition}[Supplied Plate Figure]
  \label{def:physics:phyx-mini-0777:phyxminiproblems-problemphyxmini0777-suppliedplatefigure}
  \lean{PhyXMiniProblems.ProblemPhyXMini0777.SuppliedPlateFigure}
  \uses{def:physics:phyx-mini-0777:phyxminiproblems-problemphyxmini0777-platepoint, def:physics:phyx-mini-0777:phyxminiproblems-problemphyxmini0777-appliedforcelabel, def:physics:phyx-mini-0777:phyxminiproblems-problemphyxmini0777-sidelengthlabel, def:physics:phyx-mini-0777:phyxminiproblems-problemphyxmini0777-arrowdirection}
  Literal labels and qualitative incidences transcribed from the bitmap
\end{definition}
\begin{proof}
  This is the declared model definition of Supplied Plate Figure.
\end{proof}
\begin{definition}[Rotation Axis Geometry]
  \label{def:physics:phyx-mini-0777:phyxminiproblems-problemphyxmini0777-rotationaxisgeometry}
  \lean{PhyXMiniProblems.ProblemPhyXMini0777.RotationAxisGeometry}
  \uses{def:physics:phyx-mini-0777:phyxminiproblems-problemphyxmini0777-platepoint}
  Description of the pivot axis stated in the prose
\end{definition}
\begin{proof}
  This is the declared model definition of Rotation Axis Geometry.
\end{proof}
\begin{definition}[Square Plate Torque Setup]
  \label{def:physics:phyx-mini-0777:phyxminiproblems-problemphyxmini0777-squareplatetorquesetup}
  \lean{PhyXMiniProblems.ProblemPhyXMini0777.SquarePlateTorqueSetup}
  \uses{def:physics:phyx-mini-0777:phyxminiproblems-problemphyxmini0777-lengthquantity, def:physics:phyx-mini-0777:phyxminiproblems-problemphyxmini0777-planarpositionquantity, def:physics:phyx-mini-0777:phyxminiproblems-problemphyxmini0777-planarforcequantity, def:physics:phyx-mini-0777:phyxminiproblems-problemphyxmini0777-signedtorquequantity, def:physics:phyx-mini-0777:phyxminiproblems-problemphyxmini0777-platepoint, def:physics:phyx-mini-0777:phyxminiproblems-problemphyxmini0777-appliedforcelabel, def:physics:phyx-mini-0777:phyxminiproblems-problemphyxmini0777-platematerial, def:physics:phyx-mini-0777:phyxminiproblems-problemphyxmini0777-suppliedplatefigure, def:physics:phyx-mini-0777:phyxminiproblems-problemphyxmini0777-rotationaxisgeometry}
  Independent physical quantities in the plate experiment. In particular, netTorqueAboutO is an unknown physical torque constrained only by the governing law below; it is not defined from an answer value
\end{definition}
\begin{proof}
  This is the declared model definition of Square Plate Torque Setup.
\end{proof}
\begin{definition}[Matches Problem Statement]
  \label{def:physics:phyx-mini-0777:phyxminiproblems-problemphyxmini0777-matchesproblemstatement}
  \lean{PhyXMiniProblems.ProblemPhyXMini0777.MatchesProblemStatement}
  \uses{def:physics:phyx-mini-0777:phyxminiproblems-problemphyxmini0777-lengthinmeters, def:physics:phyx-mini-0777:phyxminiproblems-problemphyxmini0777-forcemagnitudeinnewtons, def:physics:phyx-mini-0777:phyxminiproblems-problemphyxmini0777-squareplatetorquesetup}
  Numerical and categorical data stated in the problem prose
\end{definition}
\begin{proof}
  This is the declared model definition of Matches Problem Statement.
\end{proof}
\begin{definition}[Matches Primary Figure]
  \label{def:physics:phyx-mini-0777:phyxminiproblems-problemphyxmini0777-matchesprimaryfigure}
  \lean{PhyXMiniProblems.ProblemPhyXMini0777.MatchesPrimaryFigure}
  \uses{def:physics:phyx-mini-0777:phyxminiproblems-problemphyxmini0777-lengthinmeters, def:physics:phyx-mini-0777:phyxminiproblems-problemphyxmini0777-expectedapplicationpoint, def:physics:phyx-mini-0777:phyxminiproblems-problemphyxmini0777-expectedarrowdirection, def:physics:phyx-mini-0777:phyxminiproblems-problemphyxmini0777-squareplatetorquesetup}
  Primary-image evidence, including the downward F₂ arrow that conflicts with the auxiliary prose caption. No torque value or answer choice occurs here
\end{definition}
\begin{proof}
  This is the declared model definition of Matches Primary Figure.
\end{proof}
\begin{definition}[expected XCoordinate In Meters]
  \label{def:physics:phyx-mini-0777:phyxminiproblems-problemphyxmini0777-expectedxcoordinateinmeters}
  \lean{PhyXMiniProblems.ProblemPhyXMini0777.expectedXCoordinateInMeters}
  \uses{def:physics:phyx-mini-0777:phyxminiproblems-problemphyxmini0777-platepoint}
  Expected horizontal coordinate of each named point on a centered square
\end{definition}
\begin{proof}
  This is the declared model definition of expected XCoordinate In Meters.
\end{proof}
\begin{definition}[expected YCoordinate In Meters]
  \label{def:physics:phyx-mini-0777:phyxminiproblems-problemphyxmini0777-expectedycoordinateinmeters}
  \lean{PhyXMiniProblems.ProblemPhyXMini0777.expectedYCoordinateInMeters}
  \uses{def:physics:phyx-mini-0777:phyxminiproblems-problemphyxmini0777-platepoint}
  Expected vertical coordinate of each named point on a centered square
\end{definition}
\begin{proof}
  This is the declared model definition of expected YCoordinate In Meters.
\end{proof}
\begin{definition}[Matches Square Plate Geometry]
  \label{def:physics:phyx-mini-0777:phyxminiproblems-problemphyxmini0777-matchessquareplategeometry}
  \lean{PhyXMiniProblems.ProblemPhyXMini0777.MatchesSquarePlateGeometry}
  \uses{def:physics:phyx-mini-0777:phyxminiproblems-problemphyxmini0777-xaxis, def:physics:phyx-mini-0777:phyxminiproblems-problemphyxmini0777-yaxis, def:physics:phyx-mini-0777:phyxminiproblems-problemphyxmini0777-lengthinmeters, def:physics:phyx-mini-0777:phyxminiproblems-problemphyxmini0777-positioninmeters, def:physics:phyx-mini-0777:phyxminiproblems-problemphyxmini0777-squareplatetorquesetup, def:physics:phyx-mini-0777:phyxminiproblems-problemphyxmini0777-expectedxcoordinateinmeters, def:physics:phyx-mini-0777:phyxminiproblems-problemphyxmini0777-expectedycoordinateinmeters}
  Cartesian realization of the square whose center O is the pivot. This contains only geometry derived from the side length, not a torque result
\end{definition}
\begin{proof}
  This is the declared model definition of Matches Square Plate Geometry.
\end{proof}
\begin{definition}[Matches Displayed Force Geometry]
  \label{def:physics:phyx-mini-0777:phyxminiproblems-problemphyxmini0777-matchesdisplayedforcegeometry}
  \lean{PhyXMiniProblems.ProblemPhyXMini0777.MatchesDisplayedForceGeometry}
  \uses{def:physics:phyx-mini-0777:phyxminiproblems-problemphyxmini0777-xaxis, def:physics:phyx-mini-0777:phyxminiproblems-problemphyxmini0777-yaxis, def:physics:phyx-mini-0777:phyxminiproblems-problemphyxmini0777-forcevectorinnewtons, def:physics:phyx-mini-0777:phyxminiproblems-problemphyxmini0777-forcemagnitudeinnewtons, def:physics:phyx-mini-0777:phyxminiproblems-problemphyxmini0777-squareplatetorquesetup}
  The force directions shown by the arrows, written in Cartesian components. The angle is kept as an independent radian readout linked to the 45° label by MatchesPrimaryFigure
\end{definition}
\begin{proof}
  This is the declared model definition of Matches Displayed Force Geometry.
\end{proof}
\begin{definition}[Satisfies Planar Torque Law]
  \label{def:physics:phyx-mini-0777:phyxminiproblems-problemphyxmini0777-satisfiesplanartorquelaw}
  \lean{PhyXMiniProblems.ProblemPhyXMini0777.SatisfiesPlanarTorqueLaw}
  \uses{def:physics:phyx-mini-0777:phyxminiproblems-problemphyxmini0777-planarvectorreadout, def:physics:phyx-mini-0777:phyxminiproblems-problemphyxmini0777-signedscalarreadout, def:physics:phyx-mini-0777:phyxminiproblems-problemphyxmini0777-planarmomentreadout, def:physics:phyx-mini-0777:phyxminiproblems-problemphyxmini0777-appliedforcelabel, def:physics:phyx-mini-0777:phyxminiproblems-problemphyxmini0777-squareplatetorquesetup}
  The governing moment-of-force law in every coherent unit system. It relates the independent net-torque quantity to the sum of the three moments about O, without supplying any numerical result
\end{definition}
\begin{proof}
  This is the declared model definition of Satisfies Planar Torque Law.
\end{proof}
\begin{definition}[Answer Choice]
  \label{def:physics:phyx-mini-0777:phyxminiproblems-problemphyxmini0777-answerchoice}
  \lean{PhyXMiniProblems.ProblemPhyXMini0777.AnswerChoice}
  Labels of the four multiple-choice answers
\end{definition}
\begin{proof}
  This is the declared model definition of Answer Choice.
\end{proof}
\begin{definition}[displayed Torque In Newton Meters]
  \label{def:physics:phyx-mini-0777:phyxminiproblems-problemphyxmini0777-displayedtorqueinnewtonmeters}
  \lean{PhyXMiniProblems.ProblemPhyXMini0777.displayedTorqueInNewtonMeters}
  \uses{def:physics:phyx-mini-0777:phyxminiproblems-problemphyxmini0777-answerchoice}
  Torque value in newton-metres printed beside each answer label
\end{definition}
\begin{proof}
  This is the declared model definition of displayed Torque In Newton Meters.
\end{proof}
\begin{definition}[recorded Answer Choice]
  \label{def:physics:phyx-mini-0777:phyxminiproblems-problemphyxmini0777-recordedanswerchoice}
  \lean{PhyXMiniProblems.ProblemPhyXMini0777.recordedAnswerChoice}
  \uses{def:physics:phyx-mini-0777:phyxminiproblems-problemphyxmini0777-answerchoice}
  The answer label recorded by the supplied dataset metadata
\end{definition}
\begin{proof}
  This is the declared model definition of recorded Answer Choice.
\end{proof}
\begin{definition}[Rounds To Displayed Torque]
  \label{def:physics:phyx-mini-0777:phyxminiproblems-problemphyxmini0777-roundstodisplayedtorque}
  \lean{PhyXMiniProblems.ProblemPhyXMini0777.RoundsToDisplayedTorque}
  \uses{def:physics:phyx-mini-0777:phyxminiproblems-problemphyxmini0777-torqueinnewtonmeters, def:physics:phyx-mini-0777:phyxminiproblems-problemphyxmini0777-squareplatetorquesetup, def:physics:phyx-mini-0777:phyxminiproblems-problemphyxmini0777-answerchoice, def:physics:phyx-mini-0777:phyxminiproblems-problemphyxmini0777-displayedtorqueinnewtonmeters}
  The exact torque agrees with a displayed value to the hundredth-place precision used for 2.50 N m
\end{definition}
\begin{proof}
  This is the declared model definition of Rounds To Displayed Torque.
\end{proof}
% --- Archon declaration topology end ---
% --- Archon physics formalization source end ---
